\begin{helper}
Let \(D\) be a \(T_s\)-free digraph on a finite vertex set \(W\), and let \(r:W\to\mathbb N\) be any ranking function. Then the ordered graph \(\Gamma(D,r)\) from \(\noderef{DigraphOrderedGraph}\) is \(K_s\)-free.
\end{helper}
\begin{proof}
Assume for contradiction that \(\Gamma(D,r)\) contains a clique \(S\) of size \(s\). By the meaning of a \(K_s\)-clique, \(S\) has exactly \(s\) vertices and every two distinct vertices of \(S\) are adjacent in \(\Gamma(D,r)\).

First note that the ranks of the vertices of \(S\) are pairwise distinct. Indeed, if \(u,v\in S\) are distinct and \(r(u)=r(v)\), then neither \(r(u)<r(v)\) nor \(r(v)<r(u)\) holds. The definition of \(\Gamma(D,r)\) in \(\noderef{DigraphOrderedGraph}\) then gives that \(u\) and \(v\) are not adjacent, contradicting that \(S\) is a clique.

Since \(S\) is finite and the ranks on \(S\) are pairwise distinct, list its elements in increasing rank order as
\[
  v_1,\ldots,v_s,
  \qquad
  r(v_1)<\cdots<r(v_s).
\]
This list is injective and contains exactly the elements of \(S\). For every pair \(i<j\), the vertices \(v_i\) and \(v_j\) are adjacent in \(\Gamma(D,r)\). Because \(r(v_i)<r(v_j)\), the second alternative in the definition of \(\Gamma(D,r)\) is impossible, and the first alternative gives the directed edge \((v_i,v_j)\in A(D)\).

Thus the injective ordered \(s\)-tuple \(v_1,\ldots,v_s\) has every earlier-to-later arc in \(D\). This is precisely a transitive tournament of size \(s\) in the sense of \(\noderef{TransitiveTournamentFree}\), contradicting the hypothesis that \(D\) is \(T_s\)-free. Therefore no such clique \(S\) exists.
\end{proof}
